\vspace{-2mm}
\section{Introduction}
\vspace{-4mm}
\label{sec:intro}
Transformers with softmax attention \citep{vaswani2017attention}  enjoy efficient parallel training but suffer from quadratic (in sequence length) complexity, thus motivating more RNN-like models that allow for linear-time sequence modeling. Linear attention, which replaces the exponential similarity function with a simple dot product over (possibly transformed) key/query vectors, has emerged as a promising alternative to classic softmax attention \citep{katharopoulos2020transformers,performer,kasai-etal-2021-finetuning,peng2021random}. An attractive property of linear attention  is that it admits a  ``recurrent form'' in which it can be formulated as a linear RNN with 2D hidden states \citep{katharopoulos2020transformers}, thus enabling linear-time  inference. For training, linear attention  also admits a subquadratic ``chunkwise parallel form'' which divides the sequence into non-overlapping chunks and performs (serial) inter-chunk recurrent computations followed by (parallel) intra-chunk computations \citep{GAU,sun2023retentive,VQ-Transformer}, thus (partially) maintaining  parallel training. However,  existing algorithms for linear attention are not I/O aware and thus, in practice, slower than  optimized implementations of softmax attention \cite{flashattention1,flashattention2} on moderate sequence lengths.

From a performance standpoint,  linear attention has generally been found to underperform ordinary softmax attention, often by a significant margin in language modeling \cite{kasai-etal-2021-finetuning}. Recent variants of linear attention such as RetNet \citep{sun2023retentive} and TransNormerLLM \citep{qin2023scaling} obtain significant improvements by multiplying the current hidden state with a decay factor before the RNN update. However,  these works use a global, \emph{data-independent} decay factor, despite the fact that in 1D RNNs, a \emph{data-dependent} gating mechanism has been shown to be crucial for performance \citep{unreasonable-forget-gate,HGRN}. And even with the decay factor, linear attention Transformers underperform the strongest Transformer architectures when pretrained from scratch. 

This work develops a hardware-efficient algorithm for linear attention, and applies it to train a gated variant of linear attention that is competitive with softmax attention. We first discuss  aspects of optimizing ordinary linear attention on modern GPUs  and give two I/O-aware algorithms (tailored for different training settings) based on these principles (\S\ref{sec:algorithm}). Our implementation of the algorithm, called \textsc{FlashLinearAttention}, is faster than \textsc{FlashAttention-2} \citep{flashattention2} even on short (e.g., 1K) sequences. 
We then describe a gated linear attention layer with a data-dependent gating mechanism and show how \textsc{FlashLinearAttention} can be generalized to the gated case (\S\ref{sec:gla}). 
We study the resulting \emph{gated linear attention (GLA) Transformer} on moderate-scale language modeling benchmarks, where we train models with 340M/1.3B parameters on 15B/100B tokens, respectively. We find that the GLA Transformer performs favorably against a strong LLaMA architecture Transformer baseline that makes use of recent recipes \citep[Transformer++;][]{touvron2023llama} as well as recent linear-time sequence models such as RetNet \cite{sun2023retentive} and Mamba \citep{Gu2023MambaLS}. GLA Transformer is found to be particularly strong at length generalization and recall-intensive tasks among linear recurrent models. For training speed,  the GLA Transformer has significantly  higher throughput than a similarly sized Mamba model.
%
